% ============================================================================
\chapter{Conclusions and Recommendations}
\label{ch:conclusions}
% ============================================================================

This thesis set out to address a combined apparatus-and-characterisation gap in the current assessment of railway ballast degradation under drop-weight impact loading, specifically in the context of the RDSO \qtyrange{20}{65}{\mm} gradation envelope and the \qty{22.5}{\tonne} Indian Railways axle-load stress range. The programme comprised: the design, fabrication, and commissioning of a purpose-built large-scale drop-weight impact apparatus at Gati Shakti Vishwavidyalaya; the execution of a nine-configuration factorial test matrix crossing three pad conditions (UR / RUBM / PUBM) with two base conditions (FB / RB) and two impact energies (\qty{300}{\mm} / \qty{450}{\mm} drop); the characterisation of bulk degradation through the Indian Railway Fouling Index $\FIIN$ on the post-test gradations; the paired three-dimensional morphological characterisation of 45 individually tracked tracer particles through twenty-two descriptors before and after loading; and the statistical analysis of both characterisations individually and in synthesis. The present chapter states the principal conclusions of the programme, enumerates the specific novel contributions, identifies the limitations that bound the generalisability of the findings, and recommends concrete directions for follow-up research.

\section{Conclusions}
\label{sec:conclusions}

Six principal conclusions are drawn from the experimental programme and its analysis. They are stated here in the order in which they emerged during the research.
\begin{enumerate}[label=\arabic*.,leftmargin=2.2em]
  \item A large-scale drop-weight set-up was successfully developed at GSV, delivering consistent and reproducible impact energies representative of Indian Railways conditions.
  \item Ballast shows rapid initial settlement that stabilises after about fifteen blows, confirming a twenty-blow test as appropriate. Higher energy and rigid bases cause greater settlement.
  \item Fouling varies significantly with test conditions---highest under a rigid base, no pad, and high energy; lowest with a polyurethane pad, flexible base, and low energy.
  \item Settlement and fouling measure different mechanisms (rearrangement versus breakage), so both are required for a full assessment.
  \item Abrasion (Mode~A): all particles undergo gradual corner abrasion, increasing sphericity and reducing size across all configurations. Fracture (Mode~B): severe particle breakage occurs only under rigid-base conditions, owing to higher stress reflection.
  \item $\FIIN$ captures bulk fouling effects, while particle morphology reveals damage mechanisms; both are essential for complete characterisation.
\end{enumerate}

\section{Novel contributions}
\label{sec:contributions}

\subsection{Apparatus and method}
\label{sec:contrib-apparatus}

\begin{enumerate}[label=(\alph*),leftmargin=2.2em]
  \item A modified drop-weight impact apparatus for Indian ballast. A utility patent has been filed with the Indian Patent Office.
  \item A standardised protocol calibrated to the Indian Railways stress window (\qtyrange{275}{523}{\kPa}), recording per-blow peak force and five-milestone settlement---a full force-and-deformation history---versus the single post-test fines percentage of conventional Aggregate Impact Value testing.
  \item A morphological-analysis method (\cref{app:script}) that ingests paired pre/post STL meshes and exports the descriptors to a multi-sheet Excel workbook for direct comparison.
\end{enumerate}

\subsection{Scientific findings}
\label{sec:contrib-science}

\begin{enumerate}[label=(\alph*),leftmargin=2.2em]
  \item Empirical demonstration that base condition (sand cushion versus rigid concrete) governs the selection between Mode~A (diffuse abrasion) and Mode~B (catastrophic fracture) under drop-weight impact loading of ballast, using a paired per-particle 3D morphological method.
  \item A heat map of mean per-particle percentage change across fourteen descriptors for nine configurations: sign-uniformity in the size and primary-shape descriptors, heterogeneity in the shape-ratio descriptors, and sphericity and convexity as more robust diagnostics than angularity and roundness.
  \item The first published nine-configuration factorial dataset pairing large-diameter drop-weight impact loading (RDSO gradation, Indian Railways stress levels) with per-particle 3D morphology---a reference for validating discrete-element models and machine-learning degradation predictors.
\end{enumerate}

\section{Limitations of the study}
\label{sec:limitations}

The following limitations bound the generalisability of the conclusions and should be recorded in any reading of the findings.
\begin{enumerate}[label=\arabic*.,leftmargin=2.2em]
  \item \textbf{Sample size per configuration.} Five tracer particles supports the pooled Mode~A significance test and Mode~B detection, but not per-configuration ANOVA or tight Mode~B probability estimates. Expanding to 10--20 tracers is the single most effective follow-up improvement.
  \item \textbf{Single parent rock type.} Only basalt from one Vadodara quarry was tested. Parent rock affects both the Mode~A abrasion rate and the Mode~B fracture threshold, so generalisation requires extension to another RDSO-approved rock.
  \item \textbf{Surface morphology only.} The EinScan structured-light scanner cannot detect sub-surface micro-cracks, incipient fracture planes, or mineralogical alteration. Mode~A particles may carry latent damage that accelerates Mode~B fracture; X-ray CT would be required to detect this.
\end{enumerate}

\section{Recommendations for future work}
\label{sec:recommendations}

The following concrete directions for follow-up research are identified, in rough order of increasing scope. Each would usefully build on the methodological platform established by the present thesis.
\begin{enumerate}[label=\arabic*.,leftmargin=2.2em]
  \item \textbf{Expanded tracer-count study.} Repeat the present nine-configuration programme with 15--20 tracer particles per configuration to support per-configuration ANOVA and tighter estimates of Mode~B event probabilities.
  \item \textbf{Multi-rock-type extension.} Extend the programme to granite and quartzite ballast, using the same nine-configuration matrix and the same per-particle morphology pipeline, to map the parent-rock dependence of the Mode~A / Mode~B crossover.
  \item \textbf{AI/ML-based prediction of degradation behaviour.} Use the per-particle paired morphological dataset to train supervised machine-learning models. A trained surrogate would allow rapid screening of pad--base combinations without running the full impact campaign, and feature-importance analysis would identify which morphological descriptors most strongly govern degradation outcomes.
  \item \textbf{Cyclic-plus-impact combined loading.} Develop a composite loading method that applies cyclic-triaxial loading (mimicking ordinary traffic) interrupted periodically by drop-weight impact events (mimicking transition-zone hot-spots), and report the corresponding Mode~A / Mode~B evolution. This is the closest laboratory analogue to real service loading.
\end{enumerate}
